\chapter{Green's Function Method}
\renewcommand{\currentchapterpath}{chapters/gf/}

\begin{flushright}
\textit{“举一隅不以三隅反，则不复也。” ——《论语·述而》
\footnote{From the \textit{Analects}, ``Shu Er'': ``If I raise one corner and [the student] cannot return with the other three, I do not repeat myself.'' A Green's function is that one corner: the response to a unit point source. Once it is known, superposition returns the other three---the field of an arbitrary source under the same boundary and initial conditions.}}
\end{flushright}

Earlier chapters solved definite problems mainly by separation of variables. The \term{Green's function method} takes a complementary route. A Green's function---also called a \term{point-source influence function}---is the field produced by a unit point source subject to prescribed boundary and/or initial conditions. Once that response is known, linearity yields the field of an arbitrary source by superposition.

This chapter follows the classical development for Poisson's equation, the method of images, time-dependent wave and heat problems, the impulse theorem, and (optionally) generalized Green formulas for elliptic, parabolic, and hyperbolic operators~\cite{green1828,duffy,stakgold,jackson}.

\chapterinput{poisson_green}
\chapterinput{images}
\chapterinput{time_green}
\chapterinput{impulse}
\chapterinput{generalized}
\chapterinput{exercises}
